\message{ !name(CIA_2025_Esteban_Bustamante.tex)}\documentclass[conference]{IEEEtran}
\IEEEoverridecommandlockouts
% The preceding line is only needed to identify funding in the first footnote. If that is unneeded, please comment it out.
%Template version as of 6/27/2024
\usepackage[spanish]{babel}
\usepackage{cite}
\usepackage{amsmath,amssymb,amsfonts}
\usepackage{algorithmic}
\usepackage{graphicx}
\usepackage{textcomp}
\usepackage{xcolor}
\bibliographystyle{./IEEEtran}
\def\BibTeX{{\rm B\kern-.05em{\sc i\kern-.025em b}\kern-.08em
    T\kern-.1667em\lower.7ex\hbox{E}\kern-.125emX}}
\begin{document}

\message{ !name(CIA_2025_Esteban_Bustamante.tex) !offset(58) }
\begin{figure}[ht]
  \centering
  \includegraphics[width=0.3\textwidth]{figures/4.jpeg}
  \caption{\label{fig:fig4} Pipeline del Decodificador MPEG-2}
\end{figure}

En la Figura \ref{fig:fig4} se muestra el diagrama en bloques completo del diseño. Ahora describimos brevemente cada una de sus partes:
\begin{enumerate}
  \item \textbf{MPEG-2 Input:} Recibe la secuencia comprimida en formato MPEG-2, que puede ser un flujo de transporte o un flujo elemental de video.
  
  \item \textbf{Elementary Stream Strip:} Extrae únicamente el flujo elemental de video del flujo de entrada.

  \item \textbf{Video Buffer:} Almacena temporalmente el flujo de bits para asegurar una alimentación continua al decodificador.

  \item \textbf{Syntax Parser:} Analiza la sintaxis del flujo de bits para identificar cabeceras, parámetros y estructuras de control.

  \item \textbf{Variable Length Decoding (VLD):} Decodifica los símbolos comprimidos con codificación de longitud variable (como Huffman), produciendo coeficientes y vectores de movimiento.

  \item \textbf{Run-Level Decoding, Inverse Quantizing, Inverse Zigzag:}
  \begin{itemize}
    \item \textit{Run-Level Decoding:} Reconstruye la secuencia de coeficientes.
    \item \textit{Inverse Quantizing:} Restaura la escala original de los valores cuantificados.
    \item \textit{Inverse Zigzag:} Reordena los coeficientes a su formato 2D original.
  \end{itemize}

  \item \textbf{Inverse Discrete Cosine Transform (IDCT):} Convierte los coeficientes del dominio de la frecuencia al dominio espacial, generando datos de píxeles.

  \item \textbf{Motion Compensation:} Usa cuadros anteriores o futuros y vectores de movimiento para reconstruir bloques predictivos en cuadros P y B.

  \item \textbf{Frame Store:} Almacena los cuadros decodificados que servirán como referencia para la predicción por movimiento.

  \item \textbf{Chroma Resampling:} Convierte el muestreo de color de formatos como 4:2:0 a formatos como 4:2:2 o 4:4:4 para mejorar la calidad visual. En nuestro caso se usa siempre el muestreo 4:2:0.

  \item \textbf{Memory Controller / RAM:} Gestiona el acceso a la memoria externa para almacenar cuadros, referencias y buffers intermedios.

  \item \textbf{Video Sync and Timing:} Genera las señales de sincronización necesarias para mostrar correctamente los cuadros en la pantalla.

  \item \textbf{Mixer:} Mezcla el video con elementos gráficos superpuestos como menús o subtítulos (OSD).

  \item \textbf{On-Screen Display (OSD):} Agrega elementos gráficos o informativos sobre la imagen de video.

  \item \textbf{YUV to RGB:} Convierte el video del espacio de color YUV al espacio RGB, usado por la mayoría de las pantallas.

  \item \textbf{Digital Video Out:} Produce la salida final del video ya decodificado y preparado para su visualización.

  \item \textbf{Dot Clock:} Proporciona la señal de reloj que controla el ritmo de presentación de los píxeles en pantalla.
\end{enumerate}

\section{SoC donde se implementará el proyecto}
El

\label{sec:soc-donde-se}
Para poder portar nuestro diseño es necesario tener en cuenta algunos detalles:
\begin{itemize}
\item En primer lugar, el diseño original incluía el uso de un DDR2, mientras que este hace uso de un DDR2 por lo que prácticamente el controlador del mismo se encuentra 
\end{itemize}

\nocite{*}
\pagebreak
\newpage

\bibliography{bibliografia}


\message{ !name(CIA_2025_Esteban_Bustamante.tex) !offset(57) }

\end{document}

%%% Local Variables:
%%% mode: LaTeX
%%% TeX-master: t
%%% End:
